% Part:sets-functions-relations
% Chapter: size-of-sets
% Section: enumerations

\documentclass[../../../include/open-logic-section]{subfiles}

\begin{document}

\olfileid{sfr}{siz}{enm}

\olsection{گنتی وار فہرستاں تے \usetoken{S}{enumerable} سیٹ}

\begin{editorial}
  ایہ حصہ سیٹاں دیاں گنتی وار فہرستاں بارے گل کردا اے تے اوہناں نوں
  $\PosInt$ توں ہر ہدف قدر لین والے فنکشناں دے طور تے تعریف کردا اے۔
  گھٹ ریاضیاتی پس منظر والے قاریاں لئی ایہ گل ہولی ہولی سمجھائی گئی
  اے۔ اک متبادل، مختصر روپ \olref[enm-alt]{sec} وچ دتا گیا اے، جیہڑا
  گنتی وار فہرستاں نوں وکھرے ڈھنگ نال، یعنی $\Nat$ نال اک بہ اک
  مطابقتاں (یا اوہدے کسے مُڈھلے ٹکڑے نال مطابقتاں) دے طور تے تعریف کردا اے۔
\end{editorial}

\begin{explain}
اسی سیٹاں دے !!{element}s دی فہرست بنا کے پہلے وی اوہناں دیاں مثالاں
دتیاں نیں۔ آؤ ہُن عام لفظاں وچ گل کریے کہ کسے سیٹ دے !!{element}s دی
فہرست کِویں تے کدوں بن سکدی اے، بھاویں اوہ سیٹ لامتناہی ای کیوں نہ ہووے۔
\end{explain}

\begin{defn}[گنتی وار فہرست، غیر رسمی طور تے]
غیر رسمی طور تے، سیٹ~$A$ دی اک \emph{گنتی وار فہرست}، $A$ دے
!!{element}s دی اک فہرست (جیہڑی لامتناہی وی ہو سکدی اے) اے، جیہدے
وچ $A$ دا ہر !!{element} کسے متناہی مقام اُتے آوے۔ جے $A$ دی کوئی
گنتی وار فہرست اے، تاں $A$ نوں \emph{!!{enumerable}} آکھیا جاندا اے۔
\end{defn}

\begin{explain}
گنتی وار فہرستاں بارے کجھ گلاں:
\begin{enumerate}
\item اسی صرف اوہناں فہرستاں نوں گنتی وار فہرست منّنے آں جیہڑیاں
  کسے تھاں توں شروع ہوندیاں نیں، تے جیہناں وچ پہلے عنصر توں
  بغیر ہر !!{element} دے عین پہلاں اک اکلوتا !!{element} ہوندا اے۔
  دوجے لفظاں وچ، فہرست دے پہلے !!{element} تے کسے وی ہور !!{element}
  دے وچکار صرف متناہی گنتی دے عنصر ہوندے نیں۔ خاص طور تے، ایس دا
  مطلب اے کہ گنتی وار فہرست دے ہر !!{element} دا مقام متناہی اے: پہلے
  !!{element} دا مقام~$1$، دوجے دا مقام~$2$، تے ایویں اگے۔
\item اکّو سیٹ~$A$ دیاں وکھ وکھ گنتی وار فہرستاں ہو سکدیاں نیں، جیہناں
  وچ !!{element}s دا آون دا کرم وکھ ہووے: $4$، $1$، $25$، $16$،~$9$
  پہلے پنج مربع عدداں دے (سیٹ) دی اوہنی ای چنگی گنتی وار فہرست اے
  جِنّی $1$، $4$، $9$، $16$،~$25$ اے۔
\item وادھو دہراؤ والیاں فہرستاں وی گنتی وار فہرستاں ای نیں: $1$،
  $1$، $2$، $2$، $3$، $3$،~\dots{} اوسے سیٹ دی گنتی وار فہرست اے
  جیہدی $1$، $2$، $3$،~\dots{} اے۔
\item جدوں اسی کوئی خاص گنتی وار فہرست دسدے آں، تاں کرم تے دہراؤ
  \emph{واقعی} اہم ہوندے نیں: اسی مثبت صحیح عدداں دی گنتی وار فہرست
  $1$، $2$، $3$، $1$، \dots{} نال شروع کر سکدے آں، پر جدوں اوہناں نوں
  معیاری ڈھنگ نال $1$، $2$، $3$، $4$،~\dots لکھیا جاوے تاں نمونہ
  ودھ سوکھا دِسدا اے۔
\item گنتی وار فہرست دا اک مُڈھ ہونا لازمی اے: \dots، $3$، $2$، $1$
  مثبت صحیح عدداں دی گنتی وار فہرست نہیں، کیوں جے ایس دا کوئی پہلا
  !!{element} نہیں۔ ایہ گل غیر رسمی تعریف توں کِویں نکلدی اے، ویکھن
  لئی اپنے آپ نوں پچھو، ”عدد 76 فہرست دے کیہڑے مقام اُتے آندا اے؟“
\item اگلی فہرست مثبت صحیح عدداں دی گنتی وار فہرست نہیں:
  $1$، $3$، $5$، \dots، $2$، $4$، $6$، \dots\@ مسئلہ ایہ اے کہ
  جفت عدد متناہی مقاماں دی بجائے $\infty + 1$، $\infty + 2$،
  $\infty + 3$ مقاماں اُتے آندے نیں۔
\item خالی سیٹ گنتی وار فہرست بنن جوگ اے: خالی فہرست اوہدی گنتی وار فہرست اے!{}
\end{enumerate}
\end{explain}

\begin{prop}
  جے $A$ دی کوئی گنتی وار فہرست اے، تاں اوہدی دہراؤ توں بغیر وی اک
  گنتی وار فہرست اے۔
\end{prop}

\begin{proof}
  فرض کرو $A$ دی اک گنتی وار فہرست $x_1$، $x_2$، \dots{} اے، جیہدے
  وچ ہر $x_i$، $A$ دا !!{element} اے۔ اسی دہرائے ہوئے !!{element}s
  ہٹا کے فہرست دا دہراؤ ختم کر سکدے آں۔ مثال لئی، اسی ایس نوں ایہو
  جہی نویں گنتی وار فہرست بنا سکدے آں جیہدے وچ $x_i$ نوں اوس ویلے
  لکھیے جے اوہ $A$ دا !!a{element} ہووے جیہڑا $x_1$، \dots،
  $x_{i-1}$ وچ شامل نہیں؛ تے جے $x_i$ پہلے ای $x_1$، \dots،~$x_{i-1}$
  وچ آ چکیا ہووے، تاں اوہنوں فہرست توں ہٹا دیئیے۔
\end{proof}

پچھلی دلیل وکھاؤندی اے کہ گنتی وار فہرستاں تے !!{enumerable} سیٹاں
نوں چنگی طرح سمجھن تے اوہناں بارے ثبوت دین لئی سانوں ہور ٹھیک تعریف
دی لوڑ اے۔ اگلی تعریف اوہ کم کردی اے۔

\begin{defn}[گنتی وار فہرست، رسمی طور تے]
سیٹ $A \neq \emptyset$ دی اک \emph{گنتی وار فہرست} کوئی وی
!!{surjective} فنکشن $f \colon \PosInt \to A$ اے۔
\end{defn}

\begin{explain}
آؤ تسلی کر لئیے کہ رسمی تعریف تے ممکنہ طور تے لامتناہی فہرست والی
غیر رسمی تعریف اک دوجی دے برابر نیں۔ پہلے، $\PosInt$ توں سیٹ~$A$ ول
ہر !!{surjective} فنکشن، $A$ دی گنتی وار فہرست بناؤندا اے۔ ایہو جہا
فنکشن اُتلی غیر رسمی تعریف والی فہرست $f(1)$، $f(2)$، $f(3)$، \dots
متعین کردا اے۔ کیوں جے $f$، !!{surjective} اے، ایس لئی $A$ دا ہر
!!{element} لازماً $f(n)$ دی قدر اے، کسے~$n \in \PosInt$ لئی۔ ایس
لئی $A$ دا ہر !!{element} فہرست دے کسے متناہی مقام اُتے آندا اے۔
فنکشن ممکن اے !!{injective} نہ ہووے، ایس لئی فہرست وچ دہراؤ ہو سکدا
اے، پر ایہ قابلِ قبول اے (جیویں اُتے دسیا گیا اے)۔

دوجے پاسے، $A$ دے سارے !!{element}s دی گنتی وار فہرست دتی ہووے، تاں
اسی !!a{surjective} فنکشن $f\colon \PosInt \to A$ ایویں تعریف کر
سکدے آں: $f(n)$، فہرست دا $n$واں !!{element} ہووے؛ تے جے فہرست دا
کوئی $n$واں !!{element} نہ ہووے تاں اوہدا آخری !!{element} ہووے۔
ایہ صرف اوس صورت وچ !!a{surjective} فنکشن نہیں بناؤندا جدوں $A$ خالی
اے تے، ایس لئی، فہرست وی خالی اے۔ سو ہر غیر خالی فہرست !!a{surjective}
فنکشن $f\colon \PosInt \to A$ متعین کردی اے۔
\end{explain}

\begin{defn}
  \ollabel{defn:enumerable}
  سیٹ~$A$، !!{enumerable} عین اوس ویلے اے جدوں اوہ خالی ہووے یا اوہدی
  کوئی گنتی وار فہرست ہووے۔
\end{defn}

\begin{ex}
مثبت صحیح عدداں ($\PosInt$) دی گنتی وار فہرست بناؤن والا فنکشن سادہ
شناختی فنکشن $f(n) = n$ اے۔ قدرتی عدداں $\Nat$ دی گنتی وار فہرست
بناؤن والا فنکشن $g(n) = n - 1$ اے۔
\end{ex}

\begin{ex}
اگلے دتے ہوئے فنکشن $f\colon \PosInt \to \PosInt$ تے
$g \colon \PosInt \to \PosInt$:
\begin{align*}
f(n) & = 2n \text{ تے}\\
g(n) & = 2n - 1
\end{align*}
ترتیب وار جفت مثبت صحیح عدداں تے طاق مثبت صحیح عدداں دیاں گنتی وار
فہرستاں بناؤندے نیں۔ پر دوہاں وچوں کوئی وی فنکشن $\PosInt$ دی گنتی
وار فہرست نہیں، کیوں جے دوہاں وچوں کوئی وی !!{surjective} نہیں۔
\end{ex}

\begin{prob}
مثبت مربع عدداں $1$، $4$، $9$، $16$، \dots دی گنتی وار فہرست تعریف کرو۔
\end{prob}

\begin{ex}
فنکشن $f(n) = (-1)^{n} \lceil \frac{(n-1)}{2}\rceil$ (جتھے
$\lceil x \rceil$، \emph{چھت} فنکشن اے، جیہڑا $x$ نوں اوس سب توں
نیڑے صحیح عدد ول اُتے پورا کردا اے جیہڑا اوہدے توں گھٹ نہ ہووے) صحیح عدداں دے سیٹ~$\Int$ دی گنتی
وار فہرست بناؤندا اے۔ دھیان دیو کہ $f$، مثبت تے منفی صحیح عدداں دے
وچکار ”ٹپدے“ ہوئے $\Int$ دیاں قدراں پیدا کردا اے:
\[
\begin{array}{c c c c c c c c}
f(1) & f(2) & f(3) & f(4) & f(5) & f(6) & f(7) & \dots \\ \\
- \lceil \tfrac{0}{2} \rceil & \lceil \tfrac{1}{2}\rceil & - \lceil \tfrac{2}{2} \rceil & \lceil \tfrac{3}{2} \rceil & - \lceil \tfrac{4}{2} \rceil  & \lceil \tfrac{5}{2}
\rceil & - \lceil \tfrac{6}{2} \rceil & \dots \\ \\
0 & 1 & -1 & 2 & -2 & 3 & -3 & \dots
\end{array}
\]
% OLSIZ-001 ماخذ دی درستی: f(7) دی قدر -3 اے؛ انگریزی ماخذ دی آخری قطار وچ ایہ قدر رہ گئی سی۔
تسیں $f$ نوں اگلی صورتاں نال تعریف کیتا ہویا وی سمجھ سکدے او:
\[
f(n) = \begin{cases}
  0 & \text{جے $n = 1$ اے}\\
  n/2 & \text{جے $n$ جفت اے}\\
  -(n-1)/2 & \text{جے $n$ طاق تے $>1$ اے}
  \end{cases}
\]
\end{ex}

\begin{prob}
  وکھاؤ کہ جے $A$ تے $B$، !!{enumerable} نیں، تاں $A \cup B$ وی
  گنتی وار فہرست بنن جوگ اے۔ ایہ کرن لئی، فرض کرو کہ !!{surjective} فنکشن
  $f\colon \PosInt \to A$ تے $g\colon \PosInt \to B$ نیں، تے اک
  !!a{surjective} فنکشن~$h\colon \PosInt \to A \cup B$ تعریف کر کے
  ثابت کرو کہ اوہ !!{surjective} اے۔ اوہ صورتاں وی ویکھو جدوں $A$ یا
  ~$B = \emptyset$ ہووے۔
\end{prob}

\begin{prob}
  وکھاؤ کہ جے $B \subseteq A$ تے $A$، !!{enumerable} اے، تاں~$B$ وی
  گنتی وار فہرست بنن جوگ اے۔ ایہ کرن لئی، فرض کرو کہ !!a{surjective} فنکشن
  $f\colon \PosInt \to A$ اے۔ اک !!a{surjective} فنکشن
  ~$g\colon \PosInt \to B$ تعریف کر کے ثابت کرو کہ اوہ
  !!{surjective} اے۔ جے $B = \emptyset$ ہووے تاں کی ہوندا اے؟
\end{prob}

\begin{prob}
  $n$ اُتے استقرا نال وکھاؤ کہ جے $A_1$، $A_2$، \dots، $A_n$ سارے
  !!{enumerable} نیں، تاں $A_1 \cup \dots \cup A_n$ وی گنتی وار فہرست بنن جوگ
  اے۔ تسیں ایہ حقیقت منّ سکدے او کہ جے دو سیٹ $A$ تے~$B$،
  !!{enumerable} نیں، تاں~$A \cup B$ وی گنتی وار فہرست بنن جوگ اے۔
\end{prob}

سیٹ دے !!{element}s دی فہرست بناؤندیاں پہلے $1$ویں !!{element} توں
گنتی شروع کرنا شاید ودھ قدرتی لگے، پر ریاضی دان چیزاں گنن لئی قدرتی
عدد~$\Nat$ ورتنا پسند کردے نیں۔ اوہ فہرست دے $0$ویں، $1$ویں، $2$ویں
وغیرہ !!{element}s دی گل کردے نیں۔ ایس دے مطابق، اسی گنتی وار فہرست
نوں $\Nat$ توں~$A$ ول !!a{surjective} فنکشن دے طور تے تعریف کر سکدے
آں۔ بے شک، دوویں تعریفاں اک دوجی دے برابر نیں۔

\begin{prop}\ollabel{prop:enum-shift}
  !!a{surjection} $f\colon \PosInt \to A$ عین اوس ویلے موجود اے جدوں
  !!a{surjection} $g\colon \Nat \to A$ موجود اے۔
\end{prop}

\begin{proof}
  !!a{surjection} $f\colon \PosInt \to A$ دتا ہووے، تاں اسی
  $g(n) = f(n+1)$ تعریف کر سکدے آں، سارے $n \in \Nat$ لئی۔ ایہ ویکھنا
  سوکھا اے کہ $g\colon \Nat \to A$، !!{surjective} اے۔ اُلٹ پاسے،
  !!a{surjection} $g\colon \Nat \to A$ دتا ہووے تاں
  $f(n) = g(n-1)$ تعریف کرو۔
\end{proof}

ایس توں سانوں اگلا نتیجہ ملدا اے:

\begin{cor}\ollabel{cor:enum-nat}
سیٹ $A$، !!{enumerable} عین اوس ویلے اے جدوں اوہ خالی ہووے یا کوئی
!!a{surjective} فنکشن $f\colon \Nat \to A$ موجود ہووے۔
\end{cor}

اسی اُتے گل کیتی سی کہ سیٹ~$A$ دے !!{element}s دی فہرست نوں دہراؤ
توں بغیر فہرست وچ بدلیا جا سکدا اے۔ ایہ گل گنتی وار فہرستاں لئی وی
سچ اے، پر اینوں پورے رسمی ڈھنگ نال بیان تے ثابت کرنا ذرا اوکھا اے۔ ہر
فنکشن $f\colon \PosInt \to A$ دا ہر $n \in \PosInt$ لئی تعریف شدہ
ہونا لازمی اے۔ جے $A$ وچ صرف متناہی !!{element}s ہون، تاں ظاہر اے کہ
اسی $\PosInt$ دے لامتناہی !!{element}s اُتے تعریف شدہ ایہو جہا فنکشن
نہیں بنا سکدے جیہڑا $A$ دے سارے !!{element}s نوں قدر دے طور تے لے پر
اکّو قدر کدے دو واری نہ لے۔ ایس صورت وچ، یعنی جدوں دہراؤ توں بغیر
فہرست متناہی ہووے، سانوں $f$ لئی ہور دائرۂ تعریف چُننا پئے گا، جیہدے
وچ صرف متناہی !!{element}s ہون۔ دہراؤ نہ ہون دا مطلب اے کہ $f$،
!!{injective} ہووے۔ کیوں جے اوہ !!{surjective} وی اے، ایس لئی اسی کسے
متناہی سیٹ $\{1, \dots, n\}$ یا $\PosInt$ تے~$A$ دے وچکار
!!a{bijection} لبھ رہے آں۔

\begin{prop}\ollabel{prop:enum-bij}
جے $f\colon \PosInt \to A$، !!{surjective} اے (یعنی $A$ دی گنتی وار
فہرست اے)، تاں !!a{bijection} $g\colon Z \to A$ موجود اے جتھے $Z$،
یا~$\PosInt$ اے یا $\{1, \dots, n\}$ اے، کسے~$n \in \PosInt$ لئی۔
\end{prop}

\begin{proof}
  اسی فنکشن $g$ نوں بازگشتی طور تے تعریف کردے آں: $g(1) = f(1)$
  رکھو۔ جے $g(i)$ پہلے تعریف ہو چکیا اے، تاں $g(i+1)$ نوں $f(1)$،
  $f(2)$، \dots{} دی اوہ پہلی قدر رکھو جیہڑی پہلے $g(1)$، \dots،
  $g(i)$ وچ شامل نہ ہووے، جے ایہو جہی قدر موجود ہووے۔ جے $A$ دے
  ٹھیک $n$ !!{element}s ہون، تاں $g(1)$، \dots، $g(n)$ سارے تعریف
  شدہ نیں، تے ایس طرح اسی فنکشن $g\colon \{1, \dots, n\} \to A$
  تعریف کر دتا اے۔ جے $A$ دے لامتناہی !!{element}s ہون، تاں ہر $i$
  لئی $A$ دا کوئی !!a{element}، گنتی وار فہرست $f(1)$، $f(2)$،
  \dots وچ ضرور ہووے گا، جیہڑا پہلے $g(1)$، \dots، $g(i)$ وچ شامل
  نہ ہووے۔ ایس صورت وچ اسی فنکشن $g\colon \PosInt \to A$ تعریف کر
  دتا اے۔

  فنکشن $g$، !!{surjective} اے، کیوں جے $A$ دا ہر عنصر، $f(1)$،
  $f(2)$، \dots{} وچ شامل اے (کیوں جے $f$، !!{surjective} اے) تے
  ایس لئی آخرکار $g(i)$ دی قدر بنے گا، کسے~$i$ لئی۔ ایہ !!{injective}
  وی اے، کیوں جے جے کوئی $j < i$ ایہو جہا ہوندا کہ $g(j) = g(i)$،
  تاں $g(i)$ پہلے ای $g(1)$، \dots، $g(i-1)$ وچ شامل ہوندا، جیہڑا
  $g$ دی ساڈی تعریف دے اُلٹ اے۔
\end{proof}

\begin{cor}\ollabel{cor:enum-nat-bij}
سیٹ $A$، !!{enumerable} عین اوس ویلے اے جدوں اوہ خالی ہووے یا
!!a{bijection} $f\colon N \to A$ موجود ہووے، جتھے یا $N = \Nat$ اے
یا $N = \{0, \dots, n\}$ اے، کسے $n \in \Nat$ لئی۔
\end{cor}

\begin{proof}
$A$، !!{enumerable} عین اوس ویلے اے جدوں $A$ خالی ہووے یا
!!a{surjective} $f\colon \PosInt \to A$ موجود ہووے۔
\olref{prop:enum-bij} مطابق، پچھلی گل عین اوس ویلے اے جدوں اک
!!a{bijective} فنکشن~$f\colon Z \to A$ موجود ہووے، جتھے
$Z = \PosInt$ یا $Z = \{1, \dots, n\}$ اے، کسے $n \in \PosInt$
لئی۔ \olref{prop:enum-shift} دے ثبوت والی اوسے دلیل نال، اگوں ایہ عین
اوس ویلے اے جدوں اک !!a{bijection} $g\colon N \to A$ موجود ہووے،
جتھے یا $N = \Nat$ اے یا $N = \{0, \dots, n-1\}$ اے۔
\end{proof}

\begin{prob}
  \olref[sfr][siz][enm]{defn:enumerable} مطابق، سیٹ $A$ گنتی وار فہرست
  بنن جوگ عین اوس ویلے اے جدوں $A = \emptyset$ ہووے یا !!a{surjective}
  $f\colon \PosInt \to A$ موجود ہووے۔ ”!!{enumerable} سیٹ“ نوں ٹھیک
  ٹھیک ایویں تعریف کرنا وی ممکن اے: سیٹ عین اوس ویلے گنتی وار فہرست
  بنن جوگ اے جدوں !!a{injective} فنکشن $g\colon A \to \PosInt$ موجود
  ہووے۔ وکھاؤ کہ دوویں تعریفاں اک دوجی دے برابر نیں، یعنی وکھاؤ کہ
  !!a{injective} فنکشن $g\colon A \to \PosInt$ عین اوس ویلے موجود اے
  جدوں یا $A = \emptyset$ ہووے یا !!a{surjective}
  $f\colon \PosInt \to A$ موجود ہووے۔
\end{prob}

\end{document}
